% Clase 21 --- El circuito LC y la masa-resorte son el mismo problema (§2.2).
% La analogía no es adorno: la ecuación es idéntica, así que la solución también.
% La figura pone los dos sistemas uno al lado del otro y el diccionario debajo.
\begin{center}
\begin{tikzpicture}[scale=1]

  % =============== (a) masa-resorte ===============
  \begin{scope}
    \draw[figgray, line width=1.4pt] (-2.1,-0.9) -- (-2.1,1.0);
    \foreach \yy in {-0.85,-0.6,-0.35,-0.1,0.15,0.4,0.65,0.9} {
      \draw[figgray, line width=0.6pt] (-2.25,\yy) -- (-2.1,\yy+0.12);
    }
    \draw[figamber, line width=1pt, decorate,
          decoration={coil, aspect=0.55, segment length=3.4pt, amplitude=3.6pt}]
      (-2.1,0.05) -- (-0.5,0.05);
    \draw[figline, fill=figblue!25, line width=1.1pt]
      (-0.5,-0.45) rectangle (0.5,0.55);
    \node[figlbl, anchor=center] at (0,0.05) {$M$};
    \node[figlbl, figamber, anchor=south] at (-1.3,0.3) {$K$};
    \draw[figgray, line width=1pt] (-2.3,-0.9) -- (1.9,-0.9);

    \draw[figvecres, line width=1.1pt] (0.65,0.05) -- (1.55,0.05);
    \node[figlbl, figred, anchor=south] at (1.1,0.1) {$x$};

    \node[figlbl, anchor=north, align=center] at (-0.1,-1.3)
      {(a) $M\ddot x + K x = 0$};
  \end{scope}

  % =============== (b) circuito LC ===============
  \begin{scope}[xshift=6.4cm, yshift=-0.35cm]
    \draw (0,0) to[C, l=$C$, v_=$Q/C$] (0,2.0);
    \draw (0,2.0) -- (2.2,2.0);
    \draw (2.2,2.0) to[L, l=$L$] (2.2,0);
    \draw (2.2,0) -- (0,0);
    \draw[figvecres, line width=1.1pt] (0.65,2.25) -- (1.55,2.25);
    \node[figlbl, figred, anchor=south] at (1.1,2.3) {$I=\dot Q$};

    \node[figlbl, anchor=north, align=center] at (1.1,-0.5)
      {(b) $L\ddot Q + \dfrac{Q}{C} = 0$};
  \end{scope}

  \node[figlbl, anchor=north, align=center] at (3.1,-2.05)
    {$x \leftrightarrow Q$,\quad $M \leftrightarrow L$,\quad
     $K \leftrightarrow \dfrac{1}{C}$,\quad
     $\tfrac12 M\dot x^2 \leftrightarrow \tfrac12 L I^2$,\quad
     $\tfrac12 K x^2 \leftrightarrow \tfrac12 \dfrac{Q^2}{C}$};

\end{tikzpicture}\\[0.5em]
{\small\itshape Las dos ecuaciones son literalmente la misma con otros nombres,
y eso permite importar de un golpe todo lo que ya se sabe del oscilador
armónico. El diccionario tiene una lectura física: la \textbf{inductancia hace
de masa} ---es la que se resiste a que la corriente cambie, o sea la inercia del
circuito--- y el \textbf{inverso de la capacitancia hace de constante de
resorte} ---es la que tira para descargar el condensador---. La correspondencia
alcanza también a las energías: la magnética $\tfrac12 LI^2$ juega el papel de
la cinética y la eléctrica $\tfrac12 Q^2/C$ el de la potencial elástica. De ahí
que, sin resolver nada, se pueda escribir la solución
$Q(t)=Q_0\cos(\omega t+\varphi)$ y leer la frecuencia natural del diccionario:
donde el resorte tiene $\omega=\sqrt{K/M}$, el circuito tiene
$\omega = 1/\sqrt{LC}$. Verificarlo es sustituir y comprobar que
$-L\omega^2 + 1/C = 0$. Con las condiciones iniciales del problema
($Q(0)=Q_0$, $I(0)=0$) resulta $\varphi=0$.}
\end{center}
